% Transcription — fonds Alexandre Grothendieck, shelfmark 156-1, batch 2 (pages 21-26).
% Established from the facsimile archives/batches/156-1/batch-02.pdf (a copy of
% 156-1.pdf fetched through the project's relay; no cover sheet in the batch,
% so PDF page k = folder page 20+k), per the skill transcribe-grothendieck. The
% folder is twenty-six pages; this is the second and last batch, six pages.
% Batch 1 (pages 1-20) is transcribed separately and in parallel; its last
% pages (glanced at here only for context, not transcribed) introduce a
% « modèle de 1-forme » given by an order relation, locally strictly filtering
% increasing and decreasing, which « vérifie axiomes F1 : F5 ».
%
% All six leaves are mathematics in his hand and all are transcribed; nothing
% is skipped and nothing is summarised. Page 26 carries four lines only, the
% rest of the leaf blank. His own pagination runs 19-24 across the top of the
% leaves (19 overwrites an earlier figure, as does 20 on page 22); it is noted
% once per page. One date in his hand: « (7 juin) » opening page 23 — a day
% later than the 6 June heading of chapter II (folder 156-2), so this chapter
% was still being written after the next one had begun. The inventory's
% bracketed « [Chapitre] » is the archivists'; this batch carries no title.
%
% What the batch is about. Page 21: two examples of ordered sets, the first
% (L = (Q ∩ ]-∞,r]) ⊔ (Q ∩ [r,+∞]) × E, r irrational, card E ≥ 2, with maps
% p to Q and q to E ⊔ {0}) broken off and cancelled by four diagonals; the
% second, L = L_0 ⊔ ∐ L_i with L_0 below every L_i, satisfies the local
% conditions but is not filtering increasing (« n'est pas filtrante
% croissante ! »). Page 22: axiom F_6 — two segments I, J with x in ∂I ∩ ∂J
% that define the same branch at x (I_{x,y} = J_{x,y} for some y ∈ I° ∩ J°)
% lie in a common K ∈ S with x ∈ ∂K; i.e. the I° with x ∈ ∂I defining a
% given branch form an increasing filtering family (compare F3); consequence
% a) (criterion) and b) (otherwise I_{x,u} ∩ I_{x,v} = {x}); a slanted NB in
% the left margin, read only in part. Page 23 (7 June): a model of 1-form
% satisfying F_1 : F_6 is called a « réseau »; nœuds (critical or singular
% lieux), lignes (« de préférence : arêtes »), composantes, ordre or
% multiplicité of a nœud (2 for a regular lieu); « divisible » (I° ≠ ∅ for
% every segment) and « réduit » (card I = 2, all lieux nœuds, L_r = ∅).
% Page 24: a reduced réseau on L is the same as a part S of P_2(L), i.e. a
% reduced graph (« les axiomes F1) : F6) sont vides ! »); then, turning to
% regular connected réseaux, two consecutive segments I, J meeting at y with
% br_y(I) ≠ br_y(J): case I) I ∩ J = {y} gives by the recollement axiom a
% unique K ⊃ I ∪ J (« situation bien comprise »); a first case II) is
% cancelled. Pages 25-26: case II) I ∩ J ≠ {y} restarted with w ∈ I ∩ J,
% w ≠ y, reducing to the situation ∂I = ∂J = {x,y}: he shows
% br_x(I) ≠ br_x(J) by F_6, then I ∩ J = {x,y}, and the argument through the
% axiom of regular points (F5) breaks off at « on aurait br_t » on page 26.
%
% Reading hazards fixed once here. His script L is \mathcal{L}; where he
% writes a plain L (P(L), P_2(L)) it is kept plain. P, P_2 are \mathfrak{P}.
% The interior of a segment is I^{\circ} (he writes a small raised circle),
% its boundary ∂I. « br_y » (branche en y) and the sub-segments I_{x,y} are
% his. He writes « = » for « même » in running text (page 22, and apparently
% page 21); it is kept and glossed once. His underlining is set in \emph.
% The hand is fast drafting; the arcs of pages 24-25 are set as tikz-cd bent
% arrows, following the arrowheads on the page.
%
% Readings flagged and not settled: « non vide » inserted on page 21 and the
% x' < a < a'' there (presumably x' < x < x''); the insertion over « un même
% L_i » on page 21; « l'ens. des I° » on page 22; « On a donc » opening the
% page-22 consequence; most of the page-22 margin NB; the word after
% « divisible » on page 23; the index r of L_r on page 23; the first letter
% (x or y) of the page-24 margin « ∈ ∂I ∩ ∂J » (the figure suggests y); the
% page-25 margin NB on w. Variances on the page and not repaired: page 22's
% b) writes I_{x,u} ∩ I_{x,v} where J_{x,v} is expected; page 25 writes
% J = K_{x,y} for the same K_{x,y} as I, which is his point (« donc I = J »).
%
% Pass: Opus 5.5 (claude-opus-5-5), 2026-09-26 — first pass, unchecked against the pages by a human.
\documentclass[11pt,a4paper]{article}
% Le préambule commun à toutes les transcriptions du fonds Grothendieck.
%
% Il tient en une idée : **l'incertitude se compose, elle ne se résout pas.**
% Une transcription qui lisse un mot douteux a détruit la seule chose pour
% laquelle on la faisait — un lecteur doit pouvoir savoir, à chaque endroit,
% ce qui était sur le feuillet et ce qui a été ajouté. Les six macros qui
% suivent sont donc l'appareil critique, et non de la décoration.
%
% Les mêmes commandes sont reconnues par `scripts/render.mjs`, qui produit la
% vue de lecture du site. Une macro ajoutée ici doit l'être aussi là-bas,
% faute de quoi le rendu échouera bruyamment — ce qui est voulu : un rendu
% silencieusement faux est pire qu'un rendu absent.

\ProvidesPackage{grothendieck}[2026/08/08 Transcriptions du fonds Grothendieck]

\usepackage{amsmath,amssymb,amsthm}
\usepackage{mathrsfs}
\usepackage{stmaryrd}
% Les diagrammes commutatifs sont partout dans ce fonds : c'est la langue
% même de la géométrie algébrique de Grothendieck, et une transcription qui
% les rendrait en prose ne transcrirait rien.
\usepackage{tikz-cd}
% Français par défaut — c'est la langue des feuillets — mais l'anglais est
% chargé aussi : la traduction et le résumé sont des documents anglais, et la
% typographie française (espaces avant les deux-points) y serait une faute.
% Ces documents basculent par \selectlanguage{english} en tête de corps.
\usepackage[english,french]{babel}
\usepackage{fontspec}
\usepackage{geometry}
\geometry{a4paper,margin=25mm,marginparwidth=28mm}
\usepackage{marginnote}
\usepackage{xcolor}
% ulem, not soul. `soul`'s \st and \ul are built on a character-by-character
% scan that XeLaTeX's font machinery breaks: the document fails with an
% undefined control sequence rather than degrading. ulem does the same two jobs
% and is Unicode-engine safe. `normalem` stops it hijacking \emph, which the
% transcriptions use for Grothendieck's own emphasis.
\usepackage[normalem]{ulem}
\usepackage{hyperref}
\hypersetup{colorlinks=true,linkcolor=black,urlcolor=black,pdfborder={0 0 0}}

% --- Métadonnées du lot -----------------------------------------------------
% Reprises telles quelles par le script de rendu, qui les affiche en tête de
% la vue de lecture. Elles fixent ce que le fichier prétend couvrir : sans
% elles, une transcription flotte sans point d'attache dans un fonds de seize
% mille feuillets.

\newcommand{\folder}[1]{\gdef\grothFolder{#1}}
\newcommand{\batch}[1]{\gdef\grothBatch{#1}}
\newcommand{\leaves}[2]{\gdef\grothFirst{#1}\gdef\grothLast{#2}}
\newcommand{\foldertitle}[1]{\gdef\grothFoldertitle{#1}}
\gdef\grothFolder{?}\gdef\grothBatch{?}\gdef\grothFirst{?}\gdef\grothLast{?}\gdef\grothFoldertitle{}

% --- Le résumé --------------------------------------------------------------
%
% Il ouvre la lecture modernisée, sous le titre « Résumé », et il est composé
% en petit. Ce n'est pas un ornement : c'est le seul passage du document qui
% ne soit pas des mathématiques, et un lecteur doit pouvoir le reconnaître
% avant de l'avoir lu — puis le sauter s'il connaît déjà le sujet, ou s'y
% arrêter s'il ne le connaît pas. Le filet qui le suit marque la reprise du
% corps.

\newenvironment{resume}
  {\small\setlength{\parskip}{0.45em}}
  {\par\medskip\hrule\medskip}

% --- Mots-clés ---------------------------------------------------------------
%
% En anglais, en clôture du résumé de la lecture modernisée : le vocabulaire
% moderne sous lequel chercher ce que les feuillets construisent. Le manifeste
% du site (`npm run manifest`) les extrait pour étiqueter la cote entière —
% cette ligne est la source unique des tags, il n'y a pas de fichier à côté.
\newcommand{\keywords}[1]{\par\smallskip\noindent
  {\footnotesize\textsc{Keywords} --- \textit{#1}}\par}

% --- L'appareil critique ----------------------------------------------------

% Le numéro de feuillet, en marge. C'est l'ancre qui permet de régler un
% désaccord sur une formule en regardant *un* feuillet plutôt qu'une chemise
% de six cents. Le site s'en sert aussi pour tourner les pages du fac-similé
% quand on fait défiler la transcription.
\newcommand{\leaf}[1]{%
  \marginnote{\footnotesize\color{gray!70}\textsf{#1}}%
  \label{leaf:#1}\ignorespaces}

% Illisible : on ne devine pas. Un mot inventé qui se lit comme les autres est
% le pire résultat possible.
\newcommand{\ill}{\textcolor{red!60!black}{[\,\ldots\,]}}

% Lecture douteuse : le mot est donné, et signalé comme incertain.
\newcommand{\uncertain}[1]{\uline{#1}}

% Ajout de l'éditeur — une lettre manquante, un mot sous-entendu.
\newcommand{\add}[1]{\textcolor{blue!50!black}{[#1]}}

% Biffé par Grothendieck lui-même. Ce qu'il a rayé fait partie du document :
% une rature dit souvent où la pensée a changé de direction.
\newcommand{\struck}[1]{\sout{#1}}

% Note du transcripteur, sur l'état matériel du feuillet.
\newcommand{\note}[1]{\par\smallskip{\small\color{gray!60!black}\textit{[#1]}}\par\smallskip}

% Note marginale de Grothendieck, distincte du corps du texte.
\newcommand{\marginal}[1]{\par\smallskip\noindent\hspace{1em}%
  {\small\color{gray!60!black}#1}\par\smallskip}

% --- Titre ------------------------------------------------------------------

\AtBeginDocument{%
  \begin{center}
    {\large\bfseries Fonds Alexandre Grothendieck --- cote n\textsuperscript{o}~\grothFolder}\\[2pt]
    {\itshape\grothFoldertitle}\\[4pt]
    {\small feuillets \grothFirst--\grothLast\ (lot \grothBatch)}\\[2pt]
    {\footnotesize\color{gray!60!black}Transcription établie d'après le fac-similé numérisé
     par l'Université de Montpellier.\\
     Les lectures incertaines sont soulignées, les illisibles notées [\,\ldots\,],
     les ajouts entre crochets.}
  \end{center}
  \vspace{1em}}

\endinput


\folder{156-1}
\batch{2}
\pages{21}{26}
\dating{1986}
\watermark{Édition de démonstration}
\foldertitle{[Chapitre] I. Vers une géométrie des formes (topologiques) : notes manuscrites (05/06/1986).}

\begin{document}

\page{21}
\note{p. 19 de sa pagination (le chiffre surcharge un chiffre antérieur). Le premier exemple, jusqu'à « $x < y \Longleftrightarrow px < py$ et », est barré de quatre traits obliques et s'interrompt là}

\struck{\emph{Exemple} : Considérons un nb réel irrationnel $r$, et considérons l'ens}
\[
  \mathcal{L} = \underbrace{\bigl(\mathbb{Q} \cap {]{-\infty}, r]}\bigr)}_{L_1} \amalg \underbrace{\bigl(\mathbb{Q} \cap [r, +\infty]\bigr)}_{L_2} \times E
\]
\struck{où $E$ est un ens \struck{\ill{}} quelconque de card $\geq 2$. On \struck{pose} a donc}
\[
  \mathcal{L} \xrightarrow{\;p\;} \mathbb{Q} \quad (\text{surjectif}), \qquad
  \mathcal{L} \xrightarrow{\;q\;} E \amalg \{0\}, \quad q(x) = 0 \text{ si } x \in \mathcal{L}
\]
\note{l'indice de ce dernier $\mathcal{L}$ (on attend $L_1$) n'est pas visible}
\struck{et si $x, y \in \mathcal{L}$, on pose}
\[
  x < y \Longleftrightarrow px < py \ \underline{\text{et}}
\]
\note{« $px$ » : la première lettre surcharge un autre signe ; la phrase reste inachevée}

\emph{Ex} Soit $\mathcal{L}_0$ un ordonné \add{\uncertain{non vide}} \struck{filtrant} \struck{croissant et filtrant} \struck{str. filtrant.} loc. str. \struck{filtrant} croiss. et décroissant, et telle que $\exists\, x \in \mathcal{L}_0$, $\exists\, x' < a < a''$, et soit $\mathcal{L}_i$ ($i \in I$) une famille d'ens ordonnés, satisfaisant les mêmes conditions. Soit $\mathcal{L} = \mathcal{L}_0 \amalg \coprod_i \mathcal{L}_i$. On munit $\mathcal{L}$ de la relation d'ordre
\note{« non vide » est inséré au-dessus de la ligne, après un mot biffé ; « $x' < a < a''$ » : sic, on attend $x' < x < x''$}
\[
  x < y \ \text{ssi} \ \left|\ \begin{array}{l}
    \text{ou bien } x, y \text{ dans un même } \mathcal{L}_i, \text{ et } x < y \text{ dans cet } \mathcal{L}_i \\
    \text{ou bien } x \in \mathcal{L}_0,\ y \in \mathcal{L}_i \ (i \in I)
  \end{array}\right.
\]
\note{la première ligne porte des corrections : « dans \struck{$\mathcal{L}_0$ ou} », avec au-dessus une insertion lue « $= \mathcal{L}_0$ ou » d'une lecture incertaine ; « un même » est écrit « un $=$ » ; « cet » surcharge « \ill{} »}

Elle satisfait les mêmes conditions que $\mathcal{L}_0$, $\mathcal{L}_i$ mais n'est pas filtrante croissante !

\page{22}
\note{p. 20 de sa pagination, surchargeant un chiffre antérieur}

F$_6$) Soient $I, J \in S$, \add{$x \in \mathcal{L}$} tels que $x \in \partial I$, $x \in \partial J$, et que $I$, $J$ définissent la $=$ branche i.e. $\exists\, y \in I^{\circ} \cap J^{\circ}$, \add{$y \neq x$}, tel que $I_{x,y} = J_{x,y}$. Alors $\exists\, K \in S$ tel que $x \in \partial K$, $I^{\circ} \cup J^{\circ} \subset K^{\circ}$.
\note{« $x \in \mathcal{L}$ » et « $y \neq x$ » sont écrits au-dessus de la ligne ; « = » est mis pour « même », comme ailleurs}

En d'autres termes, l'\uncertain{ens.} des $I^{\circ}$, $I \in S$ tels que $x \in \partial I$, et que $I$ définisse une branche donnée, est filtrant croissant pour l'inclusion. (on comparera avec F3)
\note{« $I^{\circ}$ » est inséré au-dessus de la ligne}

\marginal{\emph{NB} On n'exige pas \ill{} que les \uncertain{intersections} \struck{\ill{}} $I \cap J$ des $I$ \uncertain{soient} dans \struck{\ill{}} $K$. Mais \struck{\ill{}} on \ill{} celles de ces \uncertain{extrémités} \ill{} \uncertain{régulières} \ill{} \struck{dans} $K$. Comme on \ill{} \uncertain{une des propriétés} \ill{} (cor. 2 p. 7)}
\note{note oblique dans la marge gauche, à une cinquantaine de degrés, très raturée ; lue en partie seulement}

On \uncertain{a donc} : Soient $I, J \in S$, $x \in \mathcal{L}$ tels que $x \in \partial I \cap \partial J$. Alors

a) Pour qu'il existe $y \in I^{\circ} \cap J^{\circ}$ tel que $I_{x,y} = J_{x,y}$, il f. et s. qu'il existe $K \in S$ tel que $x \in \partial K$, $I \cup J \subset K$.

b) Si cette condition \uncertain{n'est pas} vérifiée, alors $\exists\, u \in I^{\circ}$, $v \in J^{\circ}$ tels que
\[
  I_{x,u} \cap I_{x,v} = \{x\}.
\]
\note{sic : on attend $J_{x,v}$}

\page{23}
\note{p. 21 de sa pagination}

(7 juin) Un\struck{e} modèle de 1-forme au sens précédent (ensemble de lieux $\mathcal{L}$, avec ens de « segments » $S \subset \mathfrak{P}(L)$, et $\forall\, I \in S$, un $\partial I \in \mathfrak{P}_2(I)$, \struck{\ill{}} de façon à satisfaire les conditions (F$_1$) : (F$_6$)) est appelé un \emph{réseau}.

Les lieux critiques \add{ou singuliers} ou de branchement sont appelés aussi \emph{nœuds} du réseau. Les comp. connexes de $\mathcal{L}_r \smallsetminus \underbrace{\mathcal{L}_s}_{\text{nœuds}}$ sont appelées les \emph{lignes} du réseau (de préférence : arêtes). Composantes connexes : les \emph{composantes} du réseau.
\note{« ou singuliers » est écrit au-dessus de la ligne ; l'indice $r$ de $\mathcal{L}_r$ est d'une lecture incertaine}

On parlera de l'\emph{ordre} ou de la \emph{multiplicité} d'un nœud (plus généralement, d'un lieu du réseau --- elle est égale à 2 pour un lieu régulier).

On dira qu'un \struck{segment} réseau est \emph{divisible} \uncertain{si} \ill{} pour tt segment $I$, $I^{\circ} \neq \emptyset$. On dira qu'il est \emph{réduit} si pour tt segment $I$, $I^{\circ} = \emptyset$ i.e. $I = \partial I$ i.e. $\operatorname{card} I = 2$, \add{(ou encore si $\mathcal{L}_r = \emptyset$)}, i.e. tous ses lieux sont des nœuds.
\note{« ou encore si $\mathcal{L}_r = \emptyset$ » est inséré sous la ligne}

\page{24}
\note{p. 22 de sa pagination}

La donnée d'une structure de réseau réduit sur $\mathcal{L}$ revient à celle d'une \struck{structure} partie $S$ de $\mathfrak{P}_2(L)$ [les axiomes F$_1$) : F$_6$) sont vides !] donc une structure de \emph{graphe} réduit.

\note{un trait horizontal sépare ici deux passages}

[Je reviens : l'étude des réseaux réguliers \add{et} connexes (et non vides, i.e. 0-connexes)]

\struck{Lemme} On va étudier les cas de figure de deux segments « consécutifs »
\marginal{Rappelons que \ill{} $y$ \emph{régulier}}

\begin{tikzcd}
x \arrow[r, no head, "I"] & y \arrow[r, no head, "J"] & z
\end{tikzcd}

$y \in \mathcal{L}_r$. \marginal{$y \in \partial I \cap \partial J$ avec} $\operatorname{br}_y(I) \neq \operatorname{br}_y(J)$, Je distingue deux cas :
\note{sur la figure, le trait de $I$ porte une pointe vers $x$ ; la première lettre de la note marginale peut se lire $x$ ou $y$, la figure suggère $y$}

I) $I \cap J = \{y\}$ alors \struck{\ill{}} par axiome de recollement, $\exists !\, K \in S$, $K \supset I \cup J$, $\partial K = \{x, z\}$, et $I = K_{x,y}$, $J = K_{\struck{x}y,z}$. Situation bien comprise.

\struck{II) $I \cap J \neq \{y\}$. \struck{On} Soit \ill{}}

\struck{a) $x = z$ \quad \ill{} $= \{y, x\}$}

\begin{tikzcd}
x \arrow[r, bend left=50, "I"] & y \arrow[l, bend left=50]
\end{tikzcd}

\note{cette figure est barrée d'un trait horizontal, comme les lignes qui l'entourent}

\struck{b) $I \cap J \neq \{y, x\}$ \ill{} ($\partial \neq \{y\}$)} \struck{\ill{}}

\struck{$\exists\, t \in I \cap J$, $t \neq y$.}
\note{les lignes b) sont couvertes de hachures obliques ; le cas II) est repris à la page suivante}

\page{25}
\note{p. 23 de sa pagination}

II) $I \cap J \neq \{y\}$ \quad \struck{Deux cas}

\struck{a) \ill{}}

\begin{tikzcd}
x \arrow[r, bend left=50, "I"] & y \arrow[l, bend left=50, "J"]
\end{tikzcd}

\note{figure barrée de traits obliques}

\struck{b) $x \neq y$} Soit $w \in I \cap J$, $w \neq y$ et considérons $\underbrace{J_{y,w}}_{\in S} \subset J$, $I_{w,y}$

\begin{tikzcd}
w \arrow[r, bend left=50, "I_{w,y}"] & y \arrow[l, bend left=50, "J_{y,w}"]
\end{tikzcd}

\note{dans la marge gauche, une petite figure analogue : un cercle passant par $x$, $t$ et $y$, flèches sur l'arc supérieur vers $y$}
\marginal{\emph{NB} $w$ régulier \uncertain{sauf} si \struck{$w \in \{x, z\}$} $w = x = z$ \struck{\ill{}} i.e. si on a $x \overset{I}{\frown} z$}

\struck{Je dis que $I_{w,y} \cap$} Je vais donc étudier d'abord la situation

\begin{tikzcd}
& t \arrow[dr, bend left=20, "I"] \arrow[d, dashed, no head] & \\
x \arrow[ur, bend left=20, no head] & t \arrow[l, bend left=20, "J"] & y \arrow[l, bend left=20, no head]
\end{tikzcd}

\note{figure redessinée : un cercle, $x$ à gauche, $y$ à droite ; l'arc supérieur $I$ passe par un point $t$ et est fléché vers $y$, l'arc inférieur $J$ passe par un point $t$ et est fléché vers $x$ ; les deux $t$ sont reliés par un trait pointillé vertical}
\marginal{$\operatorname{br}_y I \neq \operatorname{br}_y J$ \quad \emph{NB} On ne va plus supposer $y$ régulier}

\struck{$I \cap J$} : $\partial I = \partial J$, $I \cap J = ?$

Je dis que $\operatorname{br}_x(I) \neq \operatorname{br}_x(J)$. Sinon, en vertu de l'axiome F$_6$, existerait \ill{} $K \in S$, tel que $x \in \partial K$, et $I = K_{x,y}$, $J = K_{x,y}$ donc $I = J$, donc $\operatorname{br}_y(I) = \operatorname{br}_y J$, absurde.

Je dis que $I \cap J = \{x, y\}$. Si en effet $t \in I \cap J$, $t \notin \{x, y\}$, alors $t \in I^{\circ} \cap J^{\circ}$, on applique ce qui précède à $I_{x,t}$, $J_{x,t}$

\page{26}
\note{p. 24 de sa pagination ; la page ne porte que ces quatre lignes}

(où $\underset{\substack{\parallel \\ \operatorname{br}_x I}}{\operatorname{br}_x I_{x,t}} \neq \underset{\substack{\parallel \\ \operatorname{br}_x (J)}}{\operatorname{br}_x J_{x,t}}$) on aurait
\[
  \operatorname{br}_t I_{x,t} \neq \operatorname{br}_t J_{x,t},
\]
donc en vertu de l'axiome des \struck{branches} pts réguliers (F$_5$) on aurait $\operatorname{br}_t$
\note{le texte s'arrête ici}

\end{document}
